\section{Introduction}
\label{sec:introduction}

The Newton systems of interior-point methods are built from the
Hessian of a self-concordant barrier, directly or through a Schur
complement or an augmented system, and can become increasingly
ill-conditioned as the requested objective accuracy improves. This is
usually analyzed for the
logarithmic barrier or for a particular primal--dual linear system. We
ask how much of it is already forced by the feasible set and the
objective. Can a different barrier, with a different central path,
improve the growth rate of the condition number of the primal barrier
Hessian? Which geometric features determine the
individual eigenvalues? The answers matter when choosing a barrier and
when interpreting solver costs that depend on the condition number.

We answer these questions for the primal barrier Hessian restricted to
the tangent space of the affine constraints (the reduced Hessian), in a
fixed Euclidean metric. Let $P$ be a compact convex feasible set with
nonempty relative interior, and let the linear objective
$\inner{c}{x}$, to be minimized, be nonconstant on the affine hull of
$P$. Let $L(g)$ be the set of feasible points whose objective value
exceeds the optimal value by at most $g$ (the primal objective gap), and let $D(g)=\diam L(g)$. For a barrier $F$, let
$\widehat H_F(g)$ be the Hessian of $F$ (not of $\mu F$) at the central
point with gap $g$, that is, at the minimizer of
$F(x)+\inner{c}{x}/\mu$ whose gap is $g$. With $\kappa$ the ratio of the largest to the smallest eigenvalue,
our main conclusion is
\begin{equation}
 \kappa\bigl(\widehat H_F(g)\bigr)
 =\Theta\!\left(\frac{D(g)^2}{g^2}\right)
 \qquad(g\downarrow0)
 \label{eq:intro-diameter-law}
\end{equation}
for every fixed self-concordant barrier $F$ with finite barrier
parameter. The constants may depend on the fixed instance, the metric,
and the barrier parameter, but not on $g$. Changing the barrier can
change the constants and the central points, but not the rate, which is
determined by the geometry of the objective sublevel sets.

The diameter law \eqref{eq:intro-diameter-law} follows from a stronger
comparison. Let $E_F$ be the unit Dikin ellipsoid of $F$ at a point,
taken in the tangent space, and let $K_g=L(g)-L(g)$ be the difference
body of the sublevel set, which depends on neither the barrier nor the
point. Suppose that, at a point
with gap $g$, the Newton decrement $\rho$ of $F+\inner{c}{x}/\mu$,
measured in the Hessian of $F$, is less than one for some $\mu>0$. Then
$E_F\subset K_g\subset2CE_F$ at this point, where $\nu$ is the barrier
parameter and
\[
 C=\frac{\nu+2\sqrt\nu+\rho}{1-\rho},\qquad 0\leq\rho<1.
\]
Since $K_g$ has diameter $2D(g)$, and the largest ball centered at the
origin inside $K_g$ has radius of order $g$, this containment gives
\eqref{eq:intro-diameter-law}.

\subsection{Main results}

\begin{enumerate}
\item \textit{Comparison at equal objective accuracy.}
\Cref{thm:difference-body,thm:diameter-law,thm:radial-spectrum} prove
the containment, the resulting Loewner-order comparison of the
Hessians of different barriers at equal gap, the diameter law, and
bounds on every eigenvalue: the $j$th eigenvalue lies between the inverse squares
of two minimax radial widths of $K_g$, up to the factor $4C^2$. The estimates
hold uniformly over barriers with bounded parameters and Newton
decrements bounded uniformly below one (\cref{cor:uniform-rates}).
On an unbounded feasible set, the pointwise bounds hold at the points of
any compact objective sublevel set (\cref{prop:localization}).
The finite barrier parameter is needed: \cref{sec:formulation-scope}
gives a self-concordant function without one for which the law fails.
\item \textit{Sharp geometric and spectral examples.}
The law gives bounded conditioning at every unique optimum of a
polytope, including degenerate vertices, and conditioning of order
$g^{-2}$ when the optimal face has positive dimension
(\cref{prop:three-regimes}). For a fixed simplex and objectives
$c_\vartheta$ whose optimum is a unique vertex for $\vartheta>0$ and an
edge at $\vartheta=0$, the condition number is
$\Theta(\min\{\vartheta^{-2},g^{-2}\})$ uniformly in $\vartheta$
(\cref{prop:simplex-plateau}). For a
trace-normalized version of a classical $3\times3$ semidefinite program,
\cref{thm:fractional-spectrum} gives the asymptotics of all four
eigenvalues of the reduced log-barrier Hessian, with exact constants.
Adding two linear constraints gives a problem whose optimality system
has the same singularity degree, two, but on which every fixed barrier
has condition number $\Theta(g^{-1})$ instead of $\Theta(g^{-3/2})$
(\cref{prop:paired-sdp}). Thus singularity degree alone does not
determine the conditioning exponent.
\item \textit{Newton right-hand sides, and what condition numbers do
not show.} For a linear program with an optimal face of dimension $f$,
every fixed barrier Hessian at the central point with gap $g$ has $f$
eigenvalues of order one and the others of order $g^{-2}$. The
orthogonal projector onto the eigenspace of the $f$ small eigenvalues (the weak eigenspace)
differs in norm from the projector onto the tangent space of the optimal
face by $O(g)$, and by $O(g^2)$ for the logarithmic barrier
(\cref{thm:lp-all-spectra}). Hence, at an exact central point, the
Newton right-hand side has a component of relative norm $O(g)$ in that
eigenspace, and $O(g^2)$ for the logarithmic barrier. An explicit $20$-self-concordant barrier
on a rectangle attains the rate $O(g)$ along a sequence of gaps
(\cref{prop:oscillatory-barrier}). For this barrier, discarding that
right-hand-side component, although its relative size tends to zero,
produces a solution whose relative Euclidean error tends to one.
For spectra in two clusters, \cref{prop:clustered-cg} bounds the
number of conjugate-gradient iterations in exact arithmetic.
\end{enumerate}

\subsection{Relation to prior work}

The geometry of self-concordant barriers is classical
\citep{nesterov1994,renegar2001,nesterov2018}; we give short proofs of
the facts we use (Dikin containment, semiboundedness, and asymmetric
containment) to fix our constants.
\citet[Proposition 2.3.2(iii), equation (2.3.9)]{nesterov1994} compares
the local Hessian norm with the symmetric chord radius at a fixed
interior point, and so already compares the Hessians of different
barriers at the \emph{same point}. We compare Hessians at generally
different points, by comparing each with the same difference body $K_g$,
which does not depend on the point.
\citet{freund2003} studies the geometry of primal--dual level sets.
More directly, \citet[Section 5.1, Fact 5.2 and Remark
5.1]{xiongfreund2024} use central-path ellipsoids to relate the diameter
and conic radius of primal--dual sublevel sets to Hessian eigenvalues,
in order to analyze rescaling for PDHG; their work is a direct
antecedent of our geometric method. We instead compare the primal
Hessian on the tangent space with the difference body $K_g$ at a
prescribed primal gap, and obtain the diameter law and the comparison
between barriers.

\citet{renegar1996} relates the conditioning of the linear equations of
barrier methods to problem conditioning and to conjugate-gradient
solution, and \citet{pena2001,pena2002} develops a primal--dual
condition-number approach and implicit-barrier methods.
\citet[Proposition 3.2]{pena2002} parameterizes the central
path by the primal objective gap, and his Section 3.3 bounds the
condition number of the Schur matrix $A\nabla^2F(x)^{-1}A^\top$ by an
inverse-square rate in the gap, under his assumptions. We therefore
claim novelty neither for the gap parameterization nor for an
inverse-square upper bound as such; for the reduced Hessian, we
add the two-sided dependence on $D(g)$ and the comparison
between barriers.

Spectral splitting of logarithmic-barrier Hessians and its numerical
consequences were studied by \citet{wright1994,wright1998};
\citet{forsgren2002} surveys interior methods, and
\citet{wrightorban2002} treat degenerate nonlinear programs. Limits of
LP central paths and their analytic properties are classical
\citep{adler1991,guler1994,halicka1999}; our limit of the log-barrier
Hessian scaled by $\mu^2$ at a unique LP optimum, including a
degenerate vertex, is a corollary of this theory.
\citet{duistermaat2001} analyzes the boundary asymptotics of the Hessian
Riemannian structure under additional smoothness assumptions; his Remark
2.3 gives oscillatory self-concordant perturbations, and qualitative
nonconvergence of central paths has earlier examples
\citep{boltepauwels2022}. Our LP contributions are the eigenvalue orders
for arbitrary finite-parameter barriers, the rates for the
small-eigenvalue eigenspace, and, for the oscillatory example, the
sharpness of the rate $O(g)$ and the calculation of the solution error.

For semidefinite programming, error bounds and facial reduction connect
singularity degree with fractional powers
\citep{sturm2000,drusvyatskiy2017,lourenco2021}. We use a
trace-normalized version of Sturm's classical $3\times3$ example and do
not claim a new source of error bounds with exponent $1/4$.
\citet{alizadeh1998} study a primal--dual Schur matrix, \citet{wei2010}
investigate instances where strict complementarity fails, and
\citet{sremac2021} relate singularity degree to primal and dual spectral
rates along a regularization path whose constraints change with the
path parameter. Our equality constraints, and hence the tangent space of
our Hessian, are fixed; \cref{sec:fractional-sdp} compares the paths and
matrices in detail.

Finally, condition numbers alone do not determine solver cost.
Cluster-sensitive conjugate-gradient analysis is classical
\citep{axelsson1994}; we give an explicit polynomial bound on the
energy-norm error in exact arithmetic, in the form needed here. Quantum
linear-system results also depend on access to the matrix,
normalization, preparation of the right-hand side, and the required
output \citep{gilyen2019,orsucci2021}, including in quantum
interior-point methods such as the quantum speedups for LP of
\citet{apersgribling2026}. We do not derive an end-to-end quantum lower
bound from the diameter law.

To the best of our knowledge, the equal-gap comparison between
arbitrary finite-parameter barriers, with bounds that depend explicitly
on the Newton decrement, and the resulting two-sided characterization by
the diameter have not previously been stated in this form. Our other
contributions are the explicit calculations listed above and their
consequences for all barriers.

\subsection{Organization}

\Cref{sec:setup,sec:geometry,sec:widths} develop the general theory,
\cref{sec:classifications,sec:lp-limits,sec:fractional-sdp} give
classifications and exact examples,
\cref{sec:lp-spectra,sec:clustered-solves} treat Newton right-hand
sides and linear solves, and \cref{sec:formulation-scope} explains the
roles of the metric, the formulation, and the finite barrier parameter.
\Cref{sec:numerics} gives reproducible numerical illustrations with
exact rational certificates of their input hypotheses; the claims they
illustrate are proved earlier.
